%macbook 
\PhanI
\setcounter{ex}{0}
%\loai{Lý thuyết}





%%\loai{Thay số tính khoảng cách}
\begin{ex}%[3P2B2-1][Loai1] 
Sóng truyền từ M đến N dọc theo phương truyền sóng với bước sóng $\lambda=120$ cm. Biết rằng sóng tại N trễ pha hơn sóng tại M là $\dfrac{\pi}{3}$. Khoảng cách $d=MN$ là
\choice
{$d=15$ cm}
{$d=24$ cm}
{$d=30$ cm}
{\True $d=20$ cm}
\loigiai{
Ta có:$\varphi = 2\pi.\frac{d}{\lambda} \Rightarrow \frac{d}{\lambda} = \frac{1}{6}\Rightarrow d = 20\ cm$.
}
\end{ex} \haidong{20}
%\loai{Tính độ lệch pha giữa hai điểm}
\begin{ex}%[3P2K3-1][Loai1]
Một nguồn sóng cơ dao động điều hoà theo phương trình $x=A\cos(3\pi t + \dfrac{\pi}{4})$ cm. Khoảng cách giữa hai điểm gần nhất trên phương truyền sóng có độ lệch pha $\dfrac{\pi}{3}$ là 0,8 m. Tốc độ truyền sóng là
\choice
{\True 7,2 m/s}
{1,6 m/s}
{4,8 m/s}
{3,2 m/s}
\loigiai{
\begin{itemize}
\item Ta có: $T = \dfrac{2\pi}{\omega} = \dfrac{2}{3}\ s$.
\item Mặt khác: $\varphi = \dfrac{2\pi.d}{\lambda}\Rightarrow \dfrac{2\pi.0,8}{\lambda} = \dfrac{\pi}{3} \Rightarrow \lambda = 4,8\ m$
$v = \dfrac{\lambda}{T} = 7,2\ m/s$.
\end{itemize}
}

\end{ex} \haidong{20}

%\loai{Dùng điều kiện kẹp - tính tần số}

\begin{ex}%[3P2K2-1][Loai1]
Một dây đàn hồi rất dài có đầu A dao động với tần số f theo phương vuông góc với sợi dây, tốc độ truyền sóng trên dây là 4,48 m/s. Xét một điểm M trên dây cách A một khoảng 28 cm, ta thấy M luôn dao động lệch pha với A một góc $\left(2k+1\right) \dfrac{\pi}{2}$ với $k=0, 1, 2...$. Biết tần số có giá trị trong khoảng từ 22 Hz đến 28 Hz. Tần số f có giá trị là
\choice
{22 Hz}
{24 Hz}
{26 Hz}
{\True 28 Hz}
\loigiai{
Độ lệch pha của hai sóng là: $\Delta \varphi = \dfrac{2\pi.28}{\lambda} = (2k + 1)\dfrac{\pi}{2}$
\begin{align*}
\Rightarrow &\dfrac{2\pi.0,28}{\dfrac{4,48}{f}} = (2k + 1)\dfrac{\pi}{2}\\
\Rightarrow &22 \leq f = (2k + 1).4 \leq 28\\
\Rightarrow &2,25 \leq k \leq 3\\
\Rightarrow &k = 3 \Rightarrow f = 28\ Hz.
\end{align*}
}
\end{ex} \haidong{20}
%\loai{Dùng điều kiện kẹp - tính tốc độ truyền sóng}
\begin{ex}%[3P2K2-1][Loai1] 
Một sóng hình sin truyền theo phương Ox từ nguồn O với tần số 30 Hz, có tốc độ truyền sóng nằm trong khoảng từ 0,5 m/s đến 0,9 m/s. Gọi A và B là hai điểm nằm trên Ox, ở cùng một phía so với O và cách nhau 5 cm. Hai phần tử môi trường tại A và B luôn dao động ngược pha với nhau. Tốc độ truyền sóng là
\choice
{80 cm/s}
{70 cm/s}
{90 cm/s}
{\True 60 cm/s}
\loigiai{
\begin{itemize}
\item Hai phần tử dao động ngược pha nhau nên 
\begin{align*}
&d = k\lambda + \dfrac{\lambda}{2} = \dfrac{2k + 1}{2}.\lambda = \dfrac{2k + 1}{2}.\dfrac{v}{f}\\ 
\Rightarrow &v = \dfrac{2d f}{2k + 1}0,5<v<0,9\\ 
\Rightarrow &0,5< \dfrac{2d f}{2k + 1}<0,9\\
\Rightarrow &1,2<k<2,5 \Rightarrow k = 2\Rightarrow \\
&v = \dfrac{2.0,05.30}{2.2 + 1} = 0,6\ m/s = 60\ cm/s
\end{align*}
\end{itemize}
}
\end{ex} \haidong{20}



\begin{ex}%[3P2K2-1][Loai1]
Một sóng cơ lan truyền trong một môi trường với tốc độ 120 cm/s, tần số của sóng thay đổi từ 10 Hz đến 15 Hz. Hai điểm cách nhau 12,5 cm luôn dao động vuông pha. Bước sóng của sóng cơ đó là
\choice
{10,5 cm}
{12 cm}
{\True 10 cm}
{8 cm}
\loigiai{
\begin{itemize}
\item Hai điểm dao động vuông pha: $d=\left( 2k+1 \right)\dfrac{\lambda }{4}=12,5\ cm$ hay $\dfrac{\left( 2k + 1 \right)v}{4f}=12,5\Rightarrow f=2,4\left( 2k+1 \right)\ Hz$. 
\item $10\ Hz\le f\le 15\ Hz \Rightarrow 10\ Hz\le 2,4\left(2k+1 \right)\le 15\ Hz\Rightarrow 1,58\le k\le 2,625\Rightarrow k=2\Rightarrow f=12\ Hz \Rightarrow \lambda =\dfrac{v}{f}=10\ cm.$
\end{itemize}
}
\end{ex} \haidong{20}
%\loai{Tính khoảng cách PQ}
\begin{ex}%[3P2K2-2][Loai1]
Một sóng ngang hình sin lan truyền trên trục Ox với biên độ không đổi và tần số bằng 4 Hz, tốc độ truyền sóng bằng 4,8 m/s. Trên Ox có hai phần tử môi trường tại P và Q luôn chuyển động cùng tốc độ nhưng ngược chiều nhau. Trong khoảng giữa PQ có hai điểm mà phần tử môi trường tại đó dao động cùng pha với phần tử môi trường tại P. Khoảng cách PQ bằng
\choice
{8 m}
{2 m}
{\True 3 m}
{4 m}
\loigiai{
\begin{itemize}
\item Ta có: $\lambda = \dfrac{v}{f} = \dfrac{4,8}{4} = 1,2\ m$\\ 
\item Lại có $v_P = -v_Q$ hay hai phần tử P và Q dao động ngược pha với nhau
$\Rightarrow PQ = k\lambda + \dfrac{\lambda}{2}$.
\item Mà giữa P và Q có 2 điểm mà phần tử môi trường tại đó dao động cùng pha với phần tử môi trường tại P nên $PQ = 2\lambda + \dfrac{\lambda}{2} = 2,5\lambda$
$\Rightarrow PQ = 2,5\lambda = 2,5.1,2 = 3$ m.
\end{itemize}
}
\end{ex} \haidong{20}

%\loai{Dùng điều kiện kẹp - tính bước sóng}
\begin{ex}%[3P2K2-1][Loai1]
Một mũi nhọn S chạm nhẹ vào mặt nước dao động điều hòa với tần số 20 Hz thì thấy hai điểm A và B trên mặt nước cùng nằm trên một phương truyền sóng cách nhau một khoảng $d=10$ cm luôn luôn dao động ngược pha với nhau. Tốc độ truyền sóng v có giá trị $0,8\ m/s \leq v \leq 1\ m/s$. Bước sóng có giá trị là
\choice
{3,5 cm}
{4,5 cm}
{\True 4 cm}
{5 cm}
\loigiai{
\begin{itemize}
\item Hai điểm AB cách nhau một khoảng d luôn dao động ngược pha nên ta có: 
\begin{align*}
\Delta \varphi = \dfrac{2\pi.d}{\lambda} = \pi + k.2\pi \Rightarrow \lambda = \dfrac{2d}{2k + 1}.
\end{align*}
\item Lại có: $v = \lambda.f = \dfrac{2d.20}{2k + 1} = \dfrac{4}{2k + 1} \ m/s$.
\item Theo bài ra: $0,8 \leq v \leq 1 \Rightarrow 1,5 \leq k \leq 2\Rightarrow k = 2\Rightarrow \lambda = 4 \ cm$.
\end{itemize}
}
\end{ex} \haidong{20}
%\loai{Tính tốc độ truyền sóng}
\begin{ex}%[3P2K2-2][Loai1]
Sóng cơ lan truyền trên sợi dây, qua hai điểm M và N cách nhau 150 cm và M sớm pha hơn N là $\dfrac{\pi}{3}+k\pi$ (k nguyên). Từ M đến N chỉ có 3 điểm vuông pha với M. Biết tần số $f = 10$ Hz. Tốc độ truyền sóng trên dây là
\choice
{100 cm/s}
{800 cm/s}
{\True 900 cm/s}
{80 m/s}
\loigiai{
Vì chỉ có 3 điểm vuông pha với M nên: \\
$\dfrac{5\pi}{2}\le \Delta \varphi <\dfrac{7\pi}{2}$ hay
$\dfrac{5\pi}{2}\le \dfrac{\pi}{3}+k\pi <\dfrac{7\pi}{2}\Rightarrow 2,2\le k<3,2\Rightarrow k=3$\\
$\Delta \varphi =\dfrac{2\pi d}{\lambda}=\dfrac{2\pi df}{v}=\dfrac{20\pi.150}{v}=\dfrac{\pi}{3}+3\pi \Rightarrow v=900\ cm/s$.
}
\end{ex} \haidong{20}
%\loai{Tính bước sóng}


%macbook 
%\loai{Đồng nhất phương trình sóng - Tìm bước sóng}
\begin{ex}%[3P2B3-1][Loai1] 
Cho một sóng ngang có phương trình sóng là $u=8\sin2\pi \left(\dfrac{t}{0,1} - \dfrac{x}{50}\right)$ mm, trong đó x tính bằng cm, t tính bằng giây. Bước sóng là
\choice
{$\lambda=0,1$ m}
{\True $\lambda=50$ cm}
{$\lambda=8$ mm}
{$\lambda=1$ m}
\loigiai{
%Ta có: $\dfrac{2\pi.x}{50} = 2\pi.\dfrac{x}{\lambda}\Rightarrow \lambda = 50\ cm$.\\
}
\end{ex} \haidong{20}
%\loai{Viết phương trình dao động tại 1 điểm - ngược chiều truyền}
\begin{ex}%[3P2B3-1][Loai1]
Sóng truyền với tốc độ 6 m/s từ điểm O đến điểm M nằm trên cùng một phương truyền sóng cách nhau 3,4 m. Coi biên độ sóng không đổi. Biết phương trình sóng tại điểm O là $u=5\cos(5\pi t+\dfrac{\pi}{6})\ cm.$ Phương trình sóng tại M là
\choice
{$u_M=5\cos(5\pi t+\dfrac{17\pi}{6})\ cm$}
{\True $u_M=5\cos(5\pi t-\dfrac{8\pi}{3})\ cm$}
{$u_M=5\cos(5\pi t+\dfrac{4\pi}{3})\ cm$}
{$u_M=5\cos(5\pi t-\dfrac{2\pi}{3})\ cm$}
\loigiai{
\begin{itemize}
\item 
Dao động tại M trễ pha hơn dao động tại O là:
$\Delta \varphi =\dfrac{2\pi d}{\lambda}=\dfrac{2\pi d}{vT}=\dfrac{\omega d}{v}=\dfrac{5\pi.3,4}{6}=\dfrac{17\pi}{6}$.
\item 
Vậy: $ u_M=5\cos \left(10\pi t+\dfrac{\pi}{6}-\dfrac{17\pi}{6}\right)=5\cos \left(10\pi t-\dfrac{8\pi}{3}\right)\ cm$ 
\end{itemize} 
}
\end{ex} \haidong{20}

%\loai{Đồng nhất phương trình sóng - Tìm tốc độ truyền sóng}
\begin{ex}%[3P2Y3-1][Loai1]
Cho một sóng ngang truyền trong một môi trường có phương trình sóng là $u=8\cos 2\pi \left( \dfrac{t}{0. 1}-\dfrac{x}{2} \right)$ mm, trong đó x tính bằng cm, t tính bằng giây. Tốc độ truyền sóng là
\choice
{\True 20 cm/s}
{20 mm/s}
{$20\pi$ cm/s}
{$10\pi$ cm/s}
\loigiai{
Ta có: $u=8\cos 2\pi \left( \dfrac{t}{0. 1}-\dfrac{x}{2} \right)=8\cos \left( 20\pi t-\pi x \right)\ mm \Rightarrow \dfrac{2\pi }{\lambda }=\pi \Rightarrow \lambda =2\ cm\Rightarrow v=2. 10=20\ cm/s$.\\
}
\end{ex} \haidong{20}
%\loai{Viết phương trình dao động tại 1 điểm - xuôi chiều truyền}






%\loai{Li độ tại một điểm trên dây}
\begin{ex}%[3P2B3-1][Loai1]
Một sóng cơ truyền dọc theo trục Ox với phương trình $u = 5\cos(8\pi t - 0,04\pi x)$ (u và x tính bằng cm, t tính bằng s). Tại thời điểm $t = 3$ s, ở điểm có $x = 25$ cm, phần tử sóng có li độ là
\choice
{$5,0$ cm}
{\True $-5,0$ cm}
{$2,5$ cm}
{$-2,5$ cm}
\loigiai{
\begin{itemize}
\item Điểm có tọa độ $x = 25$ cm có phương trình li độ dao động là: \\$u = 5\cos(8\pi t - 0,04\pi . 25) = 5\cos(8\pi t - \pi )$ cm.
\item Tại $t = 3$ s, pha dao động phần tử này là\\ ${{\phi }_{3s}}=8\pi . 3-\pi =25\pi \equiv -\pi \Rightarrow u=-5\ cm$ (hoặc ấn máy tính trực tiếp).
\end{itemize}
}

\end{ex} \haidong{20}
%\loai{Tính vận tốc tại một điểm trên dây}

%\loai{Xác định trạng thái của một điểm trên dây}


%\loai{Tính độ lệch pha giữa hai điểm}

%\loai{Quãng đường truyền sóng}


%\loai{Tốc độ dao động - tốc độ truyền sóng}
\begin{ex}%[3P2Y3-1][Loai1]
Một sóng cơ truyền trong một môi trường dọc theo trục Ox với phương trình $u=A\cos \left( 2\pi ft-\dfrac{2\pi x}{\lambda } \right)$ cm. Tốc độ dao động cực đại của các phần tử môi trường lớn gấp 4 lần tốc độ truyền sóng khi
\choice
{8$\lambda = \pi A$}
{\True 2$\lambda = \pi A$}
{6$\lambda = \pi A$}
{4$\lambda = \pi A$}
\loigiai{
\begin{itemize}
\item Ta có: $\begin{cases} v_{\max} = \omega A = 2\pi fA \\ 
v = \lambda f
\end{cases}\Rightarrow
v_{\max} = 4v \Rightarrow 2\lambda = \pi A$.
\item 
\end{itemize}
}
\end{ex} \haidong{20}





























%macbook 
%\loai{Xác định chiều truyền sóng}
\begin{ex}%[3P2K4-1][Loai1]
\immini[thm]{Sóng truyền theo chiều phương ngang đang có dạng như hình vẽ. A đang đi xuống. Phát biểu nào là đúng?
}{\begin{tikzpicture}[thick,>=stealth']
\tkzDefPoints{3.5/0/n}
\draw [red, line width = 1.2pt, domain=0:6, samples=500]%
plot (\x, {1*cos(\x*90-90)});
\draw[->](0,0)--(6.5,0);
\draw[->](0,-1.2)--(0,1.4);
\foreach \x in {0.5,1,...,6}
{\draw [dashed,thin] (\x,-1) -- (\x,1);}
\foreach \y in {-1,1}
{\draw [dashed,thin] (0,\y) -- (6,\y);}
\node[above right] at (0,1.4){$u$};
\node[above] at (6.5,0){$x$};
\node[left] at (0,0){$O$};
\draw[->] ({2.2}, {1*cos(2.2*90-90)}) --++(-90:0.6);
\draw[fill=white] ({2.2}, {1*cos(2.2*90-90)}) circle(1.2pt) node[below left]{A};
\draw[fill=white] ({4.5}, {1*cos(4.5*90-90)}) circle(1.2pt) node[below right]{B};
\end{tikzpicture}}
\choice
{Sóng truyền từ trái sang phải và B đang đi lên}
{Sóng truyền từ trái sang phải và B đang đi xuống}
{\True Sóng truyền từ phải qua trái và B đang đi lên}
{Sóng truyền từ phải qua trái và B đang đi xuống}
\end{ex} \haidong{20}


\begin{ex}%[3P2K4-1][Loai1] 
\immini[thm]{Trên mặt thoáng một chất lỏng có một nguồn phát sóng. Tại thời điểm t, hai điểm M, N trên cùng phương truyền sóng có trạng thái dao động như hình vẽ. Gọi P là trung điểm của MN. Chiều truyền sóng và trạng thái dao động của P tại thời điểm t là
}
{\begin{tikzpicture}[scale=1,thick,>=stealth']
\pgfplotsset{width=7.5cm,height=4cm,compat=1.4}
\begin{axis}[
axis lines=middle,
xmin=-0.5,xmax=6.5,xstep=1,ymin=-2,ymax=2,ystep=0.1,
grid=none, %Có các tùy chọn: minor|major|both|none
x grid style={dashed},
y grid style={dashed},
ytick={-1.5,-1,-0.5,0,0.5,1,1.5},
xtick={0.667,1.333,...,6},
xticklabels={},
yticklabels={},
xticklabel style={black,below},
xlabel style={above=3pt},
xlabel=$x(cm)$,
ylabel style={right=1pt},
ylabel={$u$}
]
\addplot[line width=1.2pt,domain=0:6,samples=200,red]{1.5*cos(deg(0.5*pi*x-pi/2))};
\draw[fill=white] (axis cs: {2*0.667}, {1.5*cos(deg(0.5*pi*2*0.667-pi/2)})circle(1.2pt)node[above right]{M};
\draw[->] (axis cs: {2*0.667}, {1.5*cos(deg(0.5*pi*2*0.667-pi/2)})--++(-90:0.5cm);
\draw[fill=white] (axis cs: {4}, {1.5*cos(deg(0.5*pi*4-pi/2)})circle(1.2pt)node[below right]{N};
\draw[->] (axis cs: {2*2}, {1.5*cos(deg(0.5*pi*4-pi/2)})--++(90:0.5cm);
\end{axis}
\end{tikzpicture}}
\choice
{Chiều từ M đến N và P đi lên}
{Chiều từ M đến N và P đi xuống}
{Chiều từ N đến M và P đi lên}
{\True Chiều từ N đến M và P đi xuống}
\loigiai{
\begin{itemize}
\item M là trung điểm của MN, từ đồ thị ta suy ra được M đi xuống thì P cũng đi xuống.
\item N lên, M xuống $\Rightarrow$ sóng truyền từ N đến M.
\end{itemize}
}
\end{ex} \haidong{20}
%\loai{Tính bước sóng}
\begin{ex}%[3P2K4-1][Loai1] 
\immini[thm]
{Một sóng cơ đang truyền theo chiều dương của trục Ox như hình vẽ. Bước sóng là
\choice
{120 cm}
{60 cm}
{30 cm}
{\True 90 cm}}
{\begin{tikzpicture}[scale=1,thick,>=stealth']
\pgfplotsset{width=7.5cm,height=4cm,compat=1.4}
\begin{axis}[
axis lines=middle,
xmin=0,xmax=6.5,xstep=1,ymin=-2,ymax=2,ystep=0.1,
grid=major, %Có các tùy chọn: minor|major|both|none
x grid style={dashed}, y grid style={dashed},
ytick={-1.5,0,1.5},
xtick={0.667,1.333,...,6},
xticklabels={},
yticklabels={-5,0,5},
xticklabel style={black,below},
xlabel style={above=1pt},
xlabel=$x(cm)$,
ylabel style={right=1pt},
ylabel={$u$}
]
\addplot[line width=1.2pt,domain=0:6,samples=200,red]{1.5*cos(deg(0.5*pi*x-pi/3))};
\draw[fill=white] (axis cs: {2*0.667}, {0})circle(1.2pt)node[below]{30};
\end{axis}
\end{tikzpicture}}
\loigiai{
\begin{itemize}
\item Từ đồ thị ta thấy mỗi khoảng trên Ox ứng với 15 cm.
\item Biên dương và biên âm cách nhau 3 ô $\sim 45\ cm = \dfrac{\lambda}{2} \Rightarrow \lambda = 90\ cm$. 
\end{itemize}
}
\end{ex} \haidong{20}
%\loai{Tính độ lệch pha giữa hai điểm}



\begin{ex}%[3P2K4-1][Loai1] 
\immini[thm]
{Một sóng ngang hình sin truyền trên một sợi dây dài. Hình vẽ bên là hình dạng của một đoạn dây tại một thời điểm xác định. Trong quá trình lan truyền sóng, hai phần tử M và N lệch nhau pha một góc là
\haicot
{$\dfrac{2\pi}{3}$ rad}
{\True $\dfrac{5\pi}{6}$ rad}
{$\dfrac{\pi}{6}$ rad}
{$\dfrac{\pi}{3}$ rad}}
{\begin{tikzpicture}[scale=1,thick,>=stealth']
\pgfplotsset{width=7.5cm,height=4cm,compat=1.4}
\begin{axis}[
axis lines=middle,
xmin=0,xmax=6.5,xstep=1,ymin=-1.8,ymax=1.8,ystep=0.1,
grid=major, %Có các tùy chọn: minor|major|both|none
x grid style={dashed},
y grid style={dashed},
ytick={-1.5,-1,-0.5,0,0.5,1,1.5},
xtick={0.33,0.66,...,6},
xticklabels={},
yticklabels={},
xticklabel style={black,below},
xlabel style={above=1pt},
xlabel=$x(cm)$,
ylabel style={right=1pt},
ylabel={$u$}
]
\addplot[line width=1.2pt,domain=0:6,samples=200,red]{1.5*cos(deg(0.5*pi*x))};
\draw[fill=white] (axis cs: 0.66, {1.5*cos(deg(0.5*pi*0.66))})circle(1.2pt)node[above right]{M};
\draw[fill=white] (axis cs: {7*0.33}, {1.5*cos(deg(0.5*pi*7*0.33))})circle(1.2pt)node[above left]{N};
\end{axis}
\end{tikzpicture}}
\loigiai{
\begin{itemize}
\item Từ đồ thị ta tính được $\lambda \sim$ 12 ô.
\item Khoảng cách từ M đến N trên Ox $\sim$ 5 ô
$\Rightarrow$ Độ lệch pha $\Delta \varphi = \dfrac{2\pi x}{\lambda}=\dfrac{2\pi.5}{12} = \dfrac{5\pi}{6}$ rad.
\end{itemize}
}
\end{ex} \haidong{20}
%\loai{Xác định trạng thái của phần tử môi trường ở thời điểm t}


%\loai{Xác định trạng thái của phần tử môi trường ở thời điểm $t+\Delta t$}
\begin{ex}%[3P2K4-1][Loai1] 
\immini[thm]
{Tại thời điểm t nào đó sóng trên sợi dây có dạng như hình vẽ. Tại thời điểm này phần tử M đang đi lên. Xác định chiều truyền sóng và vị trí của phần tử N sau đó một phần tư chu kì.
}
{\begin{tikzpicture}[scale=1,thick,>=stealth']
\pgfplotsset{width=7.5cm,height=4cm,compat=1.4}
\begin{axis}[
axis lines=middle,
xmin=-0.5,xmax=6.5,xstep=1,ymin=-1.6,ymax=1.6,ystep=0.1,
grid=major, %Có các tùy chọn: minor|major|both|none
x grid style={dashed},y grid style={dashed},
ytick={-1.5,-1,-0.5,0,0.5,1,1.5}, xtick={0.667,1.333,...,6},
xticklabels={}, yticklabels={},
xticklabel style={black,below},
xlabel style={above=3pt},
xlabel=$x(cm)$, ylabel style={right=1pt},
ylabel={$u$}
]
\addplot[line width=1.2pt,domain=0:6,samples=200,red]{1.5*cos(deg(0.5*pi*x-pi/2))};
\draw[fill=white] (axis cs: {6.5*0.667}, {1.5*cos(deg(0.5*pi*6.5*0.667-pi/2)})circle(1.2pt)node[below]{M};
\draw[fill=white] (axis cs: {3*0.667}, {1.5*cos(deg(0.5*pi*3*0.667-pi/2)})circle(1.2pt)node[below]{N};
\draw[->] (axis cs: {6.5*0.667}, {1.5*cos(deg(0.5*pi*6.5*0.667-pi/2)})--++(90:0.5cm);
\end{axis}
\end{tikzpicture}}
\choice
{Sóng truyền từ M đến N và N ở biên trên}
{Sóng truyền từ N đến M và N ở biên trên}
{\True Sóng truyền từ M đến N và N ở biên dưới}
{Sóng truyền từ N đến M và N ở biên dưới}
\loigiai{
\begin{itemize}
\item Khi M lên thì N đi xuống (sóng truyền từ M đến N).
\item Sau $\dfrac{T}{4} \Rightarrow$ N đến biên dưới.
\end{itemize}
}
\end{ex} \haidong{20}




\PhanII
\setcounter{ex}{0}

\begin{ex}
Một sóng hình sin truyền theo chiều dương của trục $O x$ với phương trình dao động của nguồn sóng (đặt tại $O$) là $u_{O}=4 \cos 20 \pi t\ cm$. Biết tốc độ truyền sóng là $20 ~cm/s$.
\choiceTF[t]
{\True Biên độ của sóng là $4 ~cm$}
{Sóng có bước sóng là $0,2 ~m$}
{Tốc độ dao động cực đại của phần tử vật chất môi trường là $20 ~cm/s$}
{\True Ở điểm $M$ (theo hướng $O x$) cách $O$ một phần tư bước sóng, phần tử môi trường dao động với phương trình là $u_{M}=4 \cos \left(20 \pi t-\dfrac{\pi}{2}\right)~cm$}
\loigiai{
\begin{itemchoice}
\itemch
\itemch $\lambda=\dfrac{v}{f}=\dfrac{20}{10}=2~cm$
\itemch $v_{\max}=A \omega=4.20 \pi=80 \pi~cm/s$
\itemch $u_{M}=4 \cos \left(20 \pi t-\dfrac{2 \pi d}{\lambda}\right)=4 \cos \left(20 \pi t-\dfrac{2 \pi \dfrac{\lambda}{4}}{\lambda}\right)=4 \cos \left(20 \pi t-\dfrac{\pi}{2}\right)~cm$
\end{itemchoice}
}
\end{ex} \haidong{20}
\begin{ex}
\immini[thm]{Một sóng cơ hình sin lan truyền trên một sợi dây đàn hồi với tần số $20 ~Hz$. Tại thời điểm t, một đoạn của sợi dây có dạng như hình bên.
phần tử ở P đang tạm dừng còn các phần tử ở M và N đang từ vị trí cân bằng của nó đi lên. Biết khoảng cách từ M đến P theo phương truyền sóng là $30 ~cm$.}{\begin{tikzpicture}[scale=1,thick,>=stealth']
\tkzDefPoints{0/0/a, 4.3/0/b}
\draw[dashed,thin] (0,0) -- (4.3,0);

\draw[->] (0.1,0) -- (0.1,0.75);
\draw[->] (4.1,0) -- (4.1,0.75);
\draw[dashed,thin] (3.13,0) -- (3.13,-1);

\def\A{1} % Biên độ
\def\T{4} % Chu kì
\def\Phi{-100} % Pha ban đầu
{\draw [red, line width = 1.2pt, domain=0:4.2, samples=500]%
plot (\x, {+(\A)*cos(360/\T*\x+\Phi)});}
\node[below right] at (0.1,0) {$M$};
\node[above] at (3.13,0) {$P$};
\node[below right] at (4.1,0) {$N$};
\end{tikzpicture}
}
\choiceTF[t]
{\True Sóng truyền từ N đến M}
{\True Tốc độ truyền sóng là $8 ~m/s$}
{\True Điểm Q là trung điểm của MP đang đi xuống}
{Sau $0,025 ~s$ thì điểm M đang đi lên}
\loigiai{
\begin{itemchoice}
\itemch Vì sóng tại M đi lên nên M và N thuộc sườn trước của sóng. Do đó sóng truyền từ M đến N
\itemch Từ đồ thị ta thấy khoảng cách từ M đến P bằng $\dfrac{3 \lambda}{4}$ $$\Rightarrow \dfrac{3 \lambda}{4}=30 \Rightarrow \lambda=40 \Rightarrow v=\lambda f=40.20=800 ~cm/s=8 ~m/s$$.
\itemch Trung điểm MN thuộc sườn sau của sóng nên nó đi xuống
\itemch $T=\dfrac{1}{f}=\dfrac{1}{20}=0,05 \Rightarrow 0,025=\dfrac{T}{2} \Rightarrow $ Tại hai thời điểm, trạng thái điểm M ngược pha nhau, ban đầu M đang đi lên nên sau đó M đang đi xuống.
\end{itemchoice}
}
\end{ex} \haidong{20}
\begin{ex}%CD
\immini[thm]{Hình bên là đồ thị li độ u (cm) - khoảng cách x (cm) tại một thời điểm của sóng ngang. Tốc độ truyền sóng là 128 cm/s.}{\begin{tikzpicture}[scale=1,>=stealth',transform shape]
%Hệ trục toạ độ
\tkzDefPoints{0/0/x', 5/0/x, 0/-1.2/y', 0/1.3/y}
\tkzDrawSegments[->,add = 0 and .05](x',x)
\tkzDrawSegments[->,add = 0 and .05](y',y)
\tkzLabelPoint[above right](x){$x\ (cm)$}
\tkzLabelPoint[above right](y){$u\ (cm)$}
\draw  [help  lines, xstep=0.5cm,ystep=0.5cm]  (x' |- y')  grid  (x |- y);
%Chia vạch trên các trục
\foreach  \y  in  {-1,1}
{\draw[thin]  (-0.1,\y)  --  (0.1,\y);}
%Chú thích trên các trục
\foreach  \y  [evaluate=\y as \z using int(6*\y)] in  {-1,1}
{\node[left=3pt] at (0,\y){$\z$};}
\node[left] at (0,0){$O$};
%Đồ thị
\def\A{1} % Biên độ
\def\T{4} % Chu kì
\def\Phi{-45} % Pha ban đầu
{\draw [red, line width = 1.2pt, domain=0:5, samples=500]%
plot (\x, {(\A)*cos(360/\T*\x+\Phi)});}
\draw[fill=white](3.5,0)circle(1pt)node[below right]{$56$};
\end{tikzpicture}}
\choiceTF[t]
{Bước sóng của sóng là $32\ cm$}
{Tần số sóng là $4\ Hz$}
{\True Tốc độ dao động cực đại của một phần tử sóng là $24\pi\ cm/s$}
{\True Quãng đường sóng truyền đi trong 5 s là $6,4\ m$}
\loigiai{
\begin{itemchoice}
\itemch $7 \text{ô}=56\Rightarrow 1\text{ô}=8\ cm$. Bước sóng: $\lambda = 8 \text{ô}=64\ cm$
\itemch Tần số sóng: $f=\dfrac{v}{\lambda}=2\ Hz$
\itemch Tốc độ cực đại của phần tử sóng: $|v|_{\max}=\omega.A=2\pi f.A=24\pi\ cm/s$
\itemch Quãng đường sóng truyền đi trong 5 s: $S=v.t=640\ cm=6,4\ m$
\end{itemchoice}
}
\end{ex} \haidong{20}
\begin{ex}%[3P2K4-1][Loai1] 
\immini[thm]{Một sóng ngang truyền trên mặt nước có tần số 10 Hz. Tại một thời điểm nào đó các phần tử mặt nước có dạng như hình vẽ. Trong đó khoảng cách từ các vị trí cân bằng của A đến vị trí cân bằng của C là 60 cm và điểm E đang đi từ vị trí cân bằng đi xuống.} 
{\begin{tikzpicture}[scale=1,thick,>=stealth']
\pgfplotsset{width=7.5cm,height=4cm,compat=1.4}
\begin{axis}[
axis lines=middle,
xmin=-0.5,xmax=6.5,xstep=1,ymin=-3,ymax=3,ystep=0.1,
grid=none, %Có các tùy chọn: minor|major|both|none
x grid style={dashed},y grid style={dashed},
ytick={-1.5,-1,-0.5,0,0.5,1,1.5}, xtick={0.667,1.333,...,6},
xticklabels={}, yticklabels={},
xticklabel style={black,below},
xlabel style={above=3pt},
xlabel=$x(cm)$, ylabel style={right=1pt},
ylabel={$u$}
]
\addplot[line width=1.2pt,domain=0:6,samples=200,red]{1.5*cos(deg(0.5*pi*x-pi/2))};
\draw[fill=white] (axis cs: {0*0.667}, {1.5*cos(deg(0.5*pi*0*0.667-pi/2)})circle(1.2pt)node[below left,fill=white,inner sep=1pt]{A};
\draw[fill=white] (axis cs: {1}, {1.5*cos(deg(0.5*pi*1-pi/2)})circle(1.2pt)node[above]{B};
\draw[fill=white] (axis cs: {2}, {1.5*cos(deg(0.5*pi*2-pi/2)})circle(1.2pt)node[below left]{C};
\draw[fill=white] (axis cs: {3}, {1.5*cos(deg(0.5*pi*3-pi/2)})circle(1.2pt)node[below]{D};
\draw[fill=white] (axis cs: {4}, {1.5*cos(deg(0.5*pi*4-pi/2)})circle(1.2pt)node[below right]{E};
\draw[fill=white] (axis cs: {5}, {1.5*cos(deg(0.5*pi*5-pi/2)})circle(1.2pt)node[above]{F};
\end{axis}
\end{tikzpicture}}
\choiceTF[t]
{\True Điểm C đang từ vị trí cân bằng đi lên}
{Sóng truyền từ E đến A}
{Bước sóng là $60\ cm$}
{\True Tốc độ truyền sóng là $12\ m/s$}	
\loigiai{
\begin{itemchoice}
\itemch 
\itemch 
\itemch Từ hình vẽ ta có $AC = \dfrac{\lambda}{2} = 60\ cm \Rightarrow \lambda = 120\ cm$ \itemch $v = \lambda f = 12\ m/s$. 
\end{itemchoice}
}
\end{ex} \haidong{20}
\PhanIII
\setcounter{ex}{0}
\begin{ex}
Một mũi nhọn $S$ chạm nhẹ vào mặt nước dao động điều hoà với tần số $25 ~Hz$. Người ta thấy rằng hai điểm $M$ và $N$ trên mặt nước cùng nằm trên phương truyền sóng cách nhau một khoảng $18 ~cm$ luôn dao động cùng pha. Biết tốc độ truyền sóng nằm trong khoảng từ $2 ~m/s$ đến $3 ~m/s$. Tốc độ truyền sóng bằng bao nhiêu m/s?\\ \shortans[oly]{2,25} 
\loigiai{
2,25
}
\end{ex} \haidong{20}
\begin{ex}
Một sóng ngang truyền trên sợi dây rất dài với tốc độ truyền sóng là $80 ~cm/s$ và tần số sóng có giá trị từ $46 ~Hz$ đến 64 Hz. Biết hai phần tử tại hai điểm trên dây cách nhau $5 ~cm$ luôn dao động ngược pha nhau. Tần số sóng trên dây bằng bao nhiêu Hz?\\ \shortans[oly]{56} 
\loigiai{
Theo giả thiết hai phần tử tại hai điểm trên dây cách nhau $5 ~cm$ luôn dao động ngược pha nhau.
$\Rightarrow \Delta \varphi=2 \pi \dfrac{d}{\lambda}=2 \pi \dfrac{d. f}{v}=(2 k+1) \pi \Rightarrow f=\dfrac{v.(2 k+1)}{2 ~d}$
$46 \leq \dfrac{v.(2 k+1)}{2 ~d} \leq 64 \Rightarrow 2,375 \leq k \leq 3,5 \Rightarrow k=3 \Rightarrow f=56 ~Hz$
}
\end{ex} \haidong{20}


\begin{ex}
Một sóng hình sin truyền theo phương $O x$ từ nguồn O với tần số $20 ~Hz$, có tốc độ truyền sóng nằm trong khoảng từ $0,6\ m/s$ đến $1 ~m/s$. Gọi A và B là hai điểm nằm trên $O x$, ở cùng một phía so với O và cách nhau $10 ~cm$. Hai phần tử môi trường tại A và B luôn dao động ngược pha với nhau. Tốc độ truyền sóng bằng bao nhiêu $m/s$?\\ \shortans[oly]{0,8} 
\loigiai{
Theo giả thiết hai phần tử môi trường tại A và B luôn dao động ngược pha với nhau
$$
\begin{aligned}
& \Rightarrow \Delta \varphi=2 \pi \dfrac{d}{\lambda}=2 \pi \dfrac{d. f}{v}=(2 k+1) \pi \Rightarrow v=\dfrac{2 df}{2 k+1} \\
& 0,6<\dfrac{2 df}{2 k+1}<1 \Rightarrow 1,5<k<2,8 \Rightarrow k=2 \Rightarrow v=0,8 ~m/s
\end{aligned}
$$
}
\end{ex} \haidong{20}
\begin{ex}
AB là một sợi dây đàn hồi căng thẳng nằm ngang, M là một điểm trên AB cách A một khoảng bằng $12,5 ~cm$. Cho A dao động điều hòa và bắt đầu đi lên từ vị trí cân bằng. Biết bước sóng là $25 ~cm$ và tần số sóng là $5 ~Hz$. Kể từ khi A bắt đầu dao động, sau bao nhiêu giây thì M lên đến điểm cao nhất?\\ \shortans[oly]{0,15} 
\loigiai{
\begin{itemize}
\item $v=\lambda f=25.5=125 ~cm/s$
\item $t=\dfrac{d}{v}+\dfrac{T}{4}=\dfrac{d}{v}+\dfrac{1}{4 f}=\dfrac{12,5}{125}+\dfrac{1}{4.5}=0,15 ~s$
\end{itemize}
}
\end{ex} \haidong{20}
\begin{ex}
Lúc $t=0$ đầu O của dây cao su căng thẳng nằm ngang bắt đầu dao động đi lên với chu kì $2 ~s$, tạo thành sóng ngang lan truyền trên dây với tốc độ $2 ~cm/s$. Điểm M trên dây cách O một khoảng $1,4 ~cm$. Thời điểm đầu tiên để M đến điểm thấp nhất là bao nhiêu giây?\\ \shortans[oly]{2,2} 
\loigiai{
$t=\dfrac{d}{v}+\dfrac{3 ~T}{4}=\dfrac{1,4}{2}+\dfrac{3.2}{4}=2,2 ~s$
}
\end{ex} \haidong{20}
\begin{ex}
\immini[thm]{Một sóng hình sin có biên độ không đổi truyền trên một sợi dây đàn hồi rất dài với tần số là $10 ~Hz$. Ở thời điểm $t$, hình dạng của đoạn dây như hình vẽ. Các vị trí cân bằng của các phần tử trên dây cùng nằm trên trục Ox. Tốc độ truyền sóng trên dây bằng bao nhiêu cm/s?}
{\begin{tikzpicture}[scale=1,thick,>=stealth']
\tkzDefPoints{0/0/o, 4.25/0/x, 0/1.5/y, 0/-1.25/y'}
\tkzDrawSegments[->](o,x y',y)
\tkzLabelPoint[left](o){$O$}
\tkzLabelPoint[above](x){$x\ (cm)$}
\tkzLabelPoint[above](y){$u\ (cm)$}
\draw[dashed,thin] (0,1) -- (0.86,1);
\draw[dashed,thin] (0,-0.6) -- (2.1,-0.6) -- (2.1,0);
\node[left] at (0,1) {$2$};
\filldraw[gray, thick] (0,1) circle (1.5pt);
\filldraw[gray, thick] (0,-0.6) circle (1.5pt);
\filldraw[red, thick] (0,0) circle (1.5pt);
\node[left] at (0,-0.6) {$-\sqrt2$};
\node[above] at (2.1,0) {$62,5$};
\filldraw[blue, thick] (2.1,0) circle (1.5pt);
\def\A{1} % Biên độ
\def\T{3.5} % Chu kì
\def\Phi{-90} % Pha ban đầu
{\draw [red, line width = 1.2pt, domain=0:3.5, samples=500]%
plot (\x, {+(\A)*cos(360/\T*\x+\Phi)});}
\end{tikzpicture}}
\shortans[oly]{1000}
\loigiai{
1000
}
\end{ex} \haidong{20}



